\newpage
\phantomsection
\label{prf:cauchy-schwarz-inequality}
\LRAProofFor{thm:cauchy-schwarz-inequality}

\begin{remark*}[Return]
\hyperref[thm:cauchy-schwarz-inequality]{Return to Theorem}
\end{remark*}

\begin{theorem*}[Cauchy--Schwarz Inequality]
For every \(x,y\in\mathbb{R}^n\),
\[
  |x\cdot y|\leq |x|\,|y|.
\]
Equality holds if and only if \(x\) and \(y\) are linearly dependent.
\end{theorem*}

\begin{proof}[Professional Standard Proof]
\LRAProofBodyStart
If \(y=0\), then both sides are \(0\), and the equality condition is immediate.
Assume \(y\neq 0\).  For every \(t\in\mathbb{R}\),
\[
  0\leq |x-ty|^2
  = |x|^2-2t(x\cdot y)+t^2|y|^2.
\]
Choose \(t=(x\cdot y)/|y|^2\).  Substitution gives
\[
  0\leq |x|^2-\frac{(x\cdot y)^2}{|y|^2}.
\]
Hence
\[
  (x\cdot y)^2\leq |x|^2|y|^2,
\]
and taking square roots gives \(|x\cdot y|\leq |x|\,|y|\).

Equality holds exactly when the chosen square has value \(0\), namely when
\[
  x-\frac{x\cdot y}{|y|^2}y=0.
\]
Thus equality implies that \(x\) is a scalar multiple of \(y\), so \(x\) and
\(y\) are linearly dependent.  Conversely, if \(x=\lambda y\), then
\[
  |x\cdot y|=|\lambda|\,|y|^2=|\lambda y|\,|y|=|x|\,|y|,
\]
so equality holds.  This also covers the case in which one vector is zero.
\end{proof}

\begin{proof}[Detailed Learning Proof]
\LRAProofBodyStart
The key observation is that squared length is always nonnegative.  For any
real number \(t\), the vector \(x-ty\) therefore satisfies
\[
  0\leq |x-ty|^2.
\]
Expanding this square with the dot product gives a quadratic expression in
\(t\):
\[
  |x-ty|^2=|x|^2-2t(x\cdot y)+t^2|y|^2.
\]
When \(y\neq 0\), the choice \(t=(x\cdot y)/|y|^2\) minimizes this quadratic.
Since the minimum is still nonnegative,
\[
  |x|^2-\frac{(x\cdot y)^2}{|y|^2}\geq 0.
\]
Multiplying by \(|y|^2\) gives
\[
  (x\cdot y)^2\leq |x|^2|y|^2.
\]
The desired inequality follows by taking square roots.

The equality case occurs precisely when the minimized square is zero.  That
means \(x-ty=0\) for the minimizing value of \(t\), so \(x=ty\).  Hence the two
vectors are linearly dependent.  If two vectors are linearly dependent, one is
a scalar multiple of the other, and direct substitution gives equality.
\end{proof}

\begin{remark*}[Proof structure]
The proof minimizes the nonnegative quadratic \(t\mapsto |x-ty|^2\).  The
inequality comes from the minimum being nonnegative; the equality condition
comes from the minimum being zero.
\end{remark*}

\begin{dependencies}
\begin{itemize}
  \item \hyperref[def:euclidean-space]{Definition: Euclidean Space}
  \item \hyperref[def:euclidean-distance-rn]{Definition: Euclidean Distance on \(\mathbb{R}^n\)}
\end{itemize}
\end{dependencies}

\clearpage
